\documentclass[11pt]{article}

\usepackage[margin=1in]{geometry}
\usepackage{graphicx}
\usepackage{amsmath}
\usepackage{amssymb}
\usepackage{caption}

\title{NeuroSynth as a Continual Learner: \\ Scaling Out to Multi-Subject Inference \\ without Forgetting}

\author{
  Priyam Das$^{1}$, Alege ALiyu$^{2}$, Chetan Gohil$^{3}$, Roshan Shetty$^{3}$, \\
  Ritoban Kundu$^{4}$, Anirban Nath$^{4}$ \\\\
  $^{1}$Independent Researcher \quad
  $^{2}$University of Minnesota \\[0.3em]
  $^{3}$AIG Labs, India \quad
  $^{4}$Indian Institute of Technology Kharagpur
}

\date{}

\begin{document}
\begin{figure}[t]
    \centering
    \includegraphics[width=\textwidth]{figures/fig3_task_b_retention.png}
    \caption{Task B Retention Across Continual Learning Phases. Task B success rate expressed as a percentage of post-Task B baseline performance, measured after each training phase. Error bars indicate one standard deviation across six seeds.}
    \label{fig:task_b_retention}
\end{figure}

\paragraph{NeuroSynth achieved improved Task C learning relative to EWC while simultaneously preserving earlier-task knowledge.}
Final Task C performance reached 9.00\% for NeuroSynth compared to 2.00\% for EWC. Although this difference did not achieve statistical significance ($p = 0.226643$), the observed moderate effect size (Cohen's $d = 0.56$) suggests meaningful performance improvements. Collectively, these results demonstrate that NeuroSynth achieved improved balance between long-term retention and continued adaptation relative to both baseline methods.

\section{Discussion}
\end{document}
